\section{Limitations}
\label{sec:limitations}

The procedural corpus is both the principal strength and the principal boundary of this study. It makes the experiment legal, controlled, and reproducible, but its seven reference pc sets and one generation process do not represent the full distribution of contemporary compositional practice. The rule generator also supplies the training targets, so a learned model can recover its regularities without demonstrating transfer to independent human-authored music. Optional composer-provided MusicXML validation is therefore important future work.

Automatic metrics simplify richer musical concepts. Pc-set coverage is recall-like and does not penalize off-set pitch classes. Interval-vector distance is computed on the generated pitch-class support and can grow when material expands beyond a small target set. Rhythmic distance uses one fixed template realization rather than the exact stochastic rhythm used to construct each target. Gesture consistency uses hand-specified feature targets rather than perceptual labels, and its normalized rest-duration feature is not bounded by one in polyphony. Serial transformation accuracy currently aliases cyclic row-order accuracy. These measures are useful diagnostics, but no single value establishes artistic coherence.

The reranking objective applies aggregate completion to every sample. On non-serial samples, this term can reward pitch classes outside a small target set. The general corpus also combines a small pc-set condition with a twelve-tone row in serial samples; the rule target then traverses the row rather than remaining inside the smaller set. Full aggregate completion can therefore conflict with interval-vector matching in both the objective and serial target construction. The same diagnostics define the reranking score and the automatic endpoints, so a favorable controlled change measures successful optimization of those diagnostics rather than independent evidence of musical quality. No sensitivity analysis of the fixed penalty weights is available. We report serial and non-serial subsets separately, but the present design does not resolve this compositional trade-off.

The primary controlled evaluation holds one trained checkpoint fixed, while the broader configuration archive uses one training seed per configuration. Multi-seed training variance is therefore unavailable. Bootstrap intervals quantify paired test-condition variability only and are not corrected for examining multiple endpoints. The rhythm and gesture fixed-default experiments collapse label variability instead of removing those token types, and the two no-constraint rows are independent-seed replications rather than distinct architectures.

The event vocabulary favors compact quantized notation. It omits tuplets, complex meters, microtonality, extended techniques, articulation detail, beaming control, dynamics in the evaluated configuration, and engraving-level layout. One instrument label is repeated across all parts of a sample, so the current interface does not model mixed ensembles. Voice-count adherence is not scored, and the exporter writes only the requested part indices; generated events assigned to higher voice indices can therefore be omitted from MusicXML. Decoding is not grammar constrained, so incomplete or out-of-order event fields can be discarded by the parser. Sequence truncation at 256 tokens may also compress longer or denser requests. In the trained representation, a time-shift token is capped at 16 quarter-note beats and \texttt{BAR} tokens do not independently advance time during decoding, so longer silent spans can be shortened. Generation can also stop before the requested span or consume the token budget. Finally, the archived music21 export path creates measures only through the last decoded event and does not pad trailing empty measures. The number of serialized measures therefore does not consistently match the 4--16-measure condition. Structural MusicXML validity should not be equated with requested-length adherence, event preservation, or publication-quality notation.

Finally, the prepared blind expert evaluation has not yet been completed. The present results support claims about implemented symbolic constraints and file export, not claims about composer preference, perceptual coherence, or practical adoption. Those outcomes are not evaluated in the present study.
